% Configuration des packages LaTeX
% Fichier: config/packages.tex

% ============================================================================
% PACKAGES DE BASE
% ============================================================================

% Encodage et polices
\usepackage[utf8]{inputenc}
\usepackage[T1]{fontenc}
\usepackage{lmodern}

% Langue
\usepackage[french]{babel}

% Mise en page
\usepackage[top=2.5cm, bottom=2.5cm, left=2.5cm, right=2.5cm]{geometry}
\usepackage{fancyhdr}
\usepackage{titlesec}
\usepackage{tocloft}

% Images et graphiques
\usepackage{graphicx}
\usepackage{float}
\usepackage{caption}
\usepackage{subcaption}
\graphicspath{{images/}}

% Couleurs
\usepackage{xcolor}
\definecolor{securityred}{RGB}{220, 53, 69}
\definecolor{securityorange}{RGB}{253, 126, 20}
\definecolor{securityyellow}{RGB}{255, 193, 7}
\definecolor{securitygreen}{RGB}{40, 167, 69}
\definecolor{securityblue}{RGB}{0, 123, 255}
\definecolor{securitypurple}{RGB}{111, 66, 193}

% Liens et références
\usepackage{hyperref}
\hypersetup{
    colorlinks=true,
    linkcolor=securityblue,
    citecolor=securitygreen,
    urlcolor=securitypurple,
    pdftitle={Rapport d'Audit de Sécurité - OWASP Juice Shop},
    pdfauthor={Auditeur de Sécurité},
    pdfsubject={Audit de sécurité de l'application web OWASP Juice Shop},
    pdfkeywords={sécurité, audit, OWASP, Juice Shop, vulnérabilités}
}

% ============================================================================
% PACKAGES POUR LE CODE ET LES PAYLOADS
% ============================================================================

% Code source
\usepackage{listings}
\lstset{
    basicstyle=\ttfamily\small,
    breaklines=true,
    frame=single,
    numbers=left,
    numberstyle=\tiny\color{gray},
    keywordstyle=\color{securityblue},
    commentstyle=\color{securitygreen},
    stringstyle=\color{securityorange},
    showstringspaces=false,
    tabsize=2
}

% Environnements pour code
\usepackage{verbatim}
\usepackage{minted}
\usemintedstyle{vs}

% Boîtes colorées
\usepackage{tcolorbox}
\tcbuselibrary{breakable, skins, theorems}

% ============================================================================
% PACKAGES POUR LES TABLEAUX ET LES MATHÉMATIQUES
% ============================================================================

% Tableaux
\usepackage{array}
\usepackage{tabularx}
\usepackage{booktabs}
\usepackage{multirow}
\usepackage{multicol}

% Mathématiques
\usepackage{amsmath}
\usepackage{amssymb}
\usepackage{amsfonts}

% ============================================================================
% PACKAGES POUR LA SÉCURITÉ ET LA DOCUMENTATION
% ============================================================================

% Packages pour la sécurité et la documentation
\usepackage{marginnote}
\usepackage{todonotes}
\usepackage{xstring}  % Pour les manipulations de chaînes (IfSubStr, StrBefore, etc.)

% Gestion des fichiers conditionnels
\usepackage{etoolbox}
\usepackage{iftex}

% Références croisées avancées
\usepackage{cleveref}

% ============================================================================
% PACKAGES SPÉCIFIQUES AU RAPPORT DE SÉCURITÉ
% ============================================================================

% Pour les diagrammes simples
\usepackage{tikz}
\usetikzlibrary{shapes, arrows, positioning}

% Pour les listes améliorées
\usepackage{enumitem}
\setlist[itemize]{leftmargin=*, label=\textbullet}
\setlist[enumerate]{leftmargin=*, label=\arabic*.}

% Pour les citations et références
\usepackage{csquotes}

% ============================================================================
% CONFIGURATIONS SPÉCIFIQUES
% ============================================================================

% Configuration des en-têtes et pieds de page
\pagestyle{fancy}
\fancyhf{}
\fancyhead[L]{\leftmark}
\fancyhead[R]{\thepage}
\renewcommand{\headrulewidth}{0.4pt}
\renewcommand{\footrulewidth}{0pt}

% Configuration des titres de sections
\titleformat{\chapter}[display]
    {\normalfont\huge\bfseries\color{securityblue}}
    {\chaptertitlename\ \thechapter}{20pt}{\Huge}

\titleformat{\section}
    {\normalfont\Large\bfseries\color{securityblue!80!black}}
    {\thesection}{1em}{}

\titleformat{\subsection}
    {\normalfont\large\bfseries\color{securityblue!60!black}}
    {\thesubsection}{1em}{}

\titleformat{\subsubsection}
    {\normalfont\normalsize\bfseries\color{securityblue!40!black}}
    {\thesubsubsection}{1em}{}

% Espacement des sections
\titlespacing*{\chapter}{0pt}{50pt}{40pt}
\titlespacing*{\section}{0pt}{3.5ex plus 1ex minus .2ex}{2.3ex plus .2ex}
\titlespacing*{\subsection}{0pt}{3.25ex plus 1ex minus .2ex}{1.5ex plus .2ex}
\titlespacing*{\subsubsection}{0pt}{3.25ex plus 1ex minus .2ex}{1ex plus .2ex}

% Configuration de la table des matières
\renewcommand{\cftchapfont}{\normalfont\bfseries\color{securityblue}}
\renewcommand{\cftsecfont}{\normalfont}
\renewcommand{\cftsubsecfont}{\normalfont\small}
\setlength{\cftbeforesecskip}{2pt}

% Configuration des légendes
\captionsetup{
    font=small,
    labelfont=bf,
    justification=centering,
    margin=1cm
}

\captionsetup[figure]{
    name=Figure,
    labelsep=colon,
    font=small,
    skip=5pt
}

\captionsetup[table]{
    name=Tableau,
    labelsep=colon,
    font=small,
    skip=5pt
}

% Configuration des sous-figures
\captionsetup[subfigure]{
    font=footnotesize,
    labelfont=bf
}

% ============================================================================
% CONFIGURATION DES ENVIRONNEMENTS DE CODE
% ============================================================================

% Style pour le code SQL
\lstdefinestyle{sqlstyle}{
    language=SQL,
    morekeywords={SELECT, FROM, WHERE, INSERT, UPDATE, DELETE, JOIN, UNION, OR, AND, LIKE},
    commentstyle=\color{securitygreen},
    keywordstyle=\color{securityblue},
    stringstyle=\color{securityorange}
}

% Style pour le code JavaScript
\lstdefinestyle{jsstyle}{
    language=JavaScript,
    morekeywords={function, var, let, const, if, else, for, while, return, alert, document},
    commentstyle=\color{securitygreen},
    keywordstyle=\color{securitypurple},
    stringstyle=\color{securityorange}
}

% Style pour le code HTML
\lstdefinestyle{htmlstyle}{
    language=HTML,
    morekeywords={div, span, script, img, a, href, src, onclick},
    commentstyle=\color{securitygreen},
    keywordstyle=\color{securityblue},
    stringstyle=\color{securityorange}
}

% ============================================================================
% CONFIGURATION DES COULEURS POUR LA SÉCURITÉ
% ============================================================================

% Niveaux de sévérité
\colorlet{criticalcolor}{securityred}
\colorlet{highcolor}{securityorange}
\colorlet{mediumcolor}{securityyellow}
\colorlet{lowcolor}{securitygreen}
\colorlet{infocolor}{securityblue}

% ============================================================================
% CONFIGURATION DES EXTENSIONS D'IMAGES
% ============================================================================

\DeclareGraphicsExtensions{.pdf,.png,.jpg,.jpeg,.gif}

% ============================================================================
% COMMANDES UTILITAIRES GLOBALES
% ============================================================================

% Commande pour les références OWASP
\newcommand{\owasp}[1]{\textbf{OWASP #1}}

% Commande pour les niveaux de difficulté
\newcommand{\difficulty}[1]{\textbf{#1}}

% Commande pour les catégories de vulnérabilités
\newcommand{\vulncategory}[1]{\textcolor{securityblue}{\textbf{#1}}}

% Fin du fichier packages.tex